\documentclass[11pt]{article}

\usepackage[margin=1in]{geometry}
\usepackage{booktabs}
\usepackage{amsmath,amssymb,amsthm}
\usepackage{natbib}
\usepackage{hyperref}

\newtheorem{theorem}{Theorem}
\newtheorem{proposition}{Proposition}
\newtheorem{corollary}{Corollary}

\newcommand{\Y}{\mathcal{Y}}
\newcommand{\Lang}{\mathcal{L}}
\newcommand{\E}{\mathbb{E}}
\newcommand{\ptheta}{p_{\theta}}
\newcommand{\Z}{Z}

\title{Auditing CRF Posteriors with Regular-Language Event Mass}
\author{Anonymous Authors}
\date{}

\begin{document}
\maketitle

\begin{abstract}
We formalize and study posterior regular-language event mass for conditional
random fields as an audit statistic under the original model posterior.
Known product-automaton marginal inference is the computational neighborhood;
the contribution is not a new CRF--DFA inference algorithm.  Instead, the paper
separates decoded legality from posterior consistency, studies event loss as a
computable training signal, and reports empirical boundaries where event mass is
useful, saturated, weaker than generic uncertainty, or harmful when the rule is
misleading.
\end{abstract}

\section{Introduction}

Regular constraints are common in structured prediction: a BIO tag sequence
should be legal, an extracted date should match a format, and a product code may
need to satisfy a finite-state rule.  A constrained decoder can enforce such a
rule at prediction time, but the legality of the decoded output does not answer
how much probability the original model assigned to legal structures.

This paper studies that latter question.  For a CRF posterior
\(\ptheta(y \mid x)\) and a regular language \(\Lang\), we report
\(\ptheta(\Lang \mid x)\), the posterior mass assigned by the original CRF to
the regular-language event.  The denominator remains the original CRF
normalizer over all length-\(T\) label sequences; it is not replaced by a
constrained-support model.

Known product-automaton marginal inference is the computational neighborhood.
Our focus is original-posterior regular-language event mass as an auditable
semantic object, separated from constrained decoding and constrained-support
normalization.  This distinction matters when a legal decoded output coexists
with low posterior mass on the legal event.

The paper makes four conservative claims.  First, posterior regular-language
event mass is a well-defined finite CRF posterior statistic and can be computed
exactly by standard product-state transfer under explicit finite-context local
factorization assumptions.  Second, hard-constrained decoding and posterior
consistency answer different questions.  Third, the event loss
\(-\log \ptheta(\Lang \mid x)\) is a computable training signal, but not an
accuracy theorem.  Fourth, empirical diagnostics show both utility and
boundaries: event risk can rank rule-specific errors, generic uncertainty can
be stronger, event mass can saturate, and misleading rules can harm task
metrics.

\paragraph{Contributions.}
The contribution is not a replacement for WFST, constrained decoding,
regular-constrained CRFs, posterior regularization, or semantic loss.  The
contribution is a posterior-semantics and audit protocol for reporting the
original CRF posterior's mass on regular-language events, with training and
diagnostic experiments that state their limitations.

\subsection{What This Object Answers}

The key distinction is the denominator and the intervention point.  Table
\ref{tab:intro:objects} summarizes the neighboring objects that otherwise look
similar in notation.

\begin{table}[t]
\centering
\small
\caption{Nearby structured-prediction objects answer different questions.}
\label{tab:intro:objects}
\begin{tabular}{llll}
\toprule
Object & Procedure & Space & Question \\
\midrule
CRF posterior &
\(\ptheta(y\mid x)\) &
all \(y\in\Y^T\) &
what distribution did the model assign? \\
Constrained decoding &
\(\arg\max_{y\in\Lang} S_\theta(x,y)\) &
legal outputs &
what is the best legal output? \\
Constrained-support CRF &
\(\ptheta(y\mid x,y\in\Lang)\) &
legal outputs &
what posterior results after support restriction? \\
This audit &
\(\ptheta(\Lang\mid x)=\Z_{\theta,\Lang}(x)/\Z_\theta(x)\) &
all outputs in denominator &
how much original posterior mass is legal? \\
\bottomrule
\end{tabular}
\end{table}

This distinction is useful only if the scalar is reported with its limitations.
If a rule is already saturated, event mass may add little diagnostic
information.  If a rule is weakly related to task correctness, optimizing event
loss can raise event mass while degrading task metrics.  If the goal is only to
return a legal output, a constrained decoder may be the right tool.  The point
of the audit is to make these cases visible rather than to replace those
interventions.

\subsection{Closest-Prior Objection}

A natural objection is that \(\ptheta(\Lang\mid x)\) is just standard marginal
inference on a product automaton with a different name.  We agree about the
computation.  The product-state dynamic program is not the novelty claim.  The
paper's claim is that the resulting original-posterior event probability is a
useful object to report and study: it can be low when constrained Viterbi
returns a legal output, can define a semantic-loss-like training signal, and can
support diagnostics with explicit empirical boundaries.  This is why the
experiments compare posterior event mass, unconstrained decoded legality,
constrained-product decoded legality, task metrics, and uncertainty baselines
as separate quantities.

\section{Related Work}

\paragraph{CRFs and automata-based inference.}
Conditional random fields provide the base conditional sequence model
\citep{lafferty2001crf}.  Finite-state methods and weighted finite-state
transducers are a close computational neighbor for constrained sequence
decoding and automata composition \citep{mohri2002wfst}.  We do not claim that
CRF--automaton product inference is new.  We use this known neighborhood to
report a different quantity: event mass under the original posterior.

\paragraph{Constrained decoding and regular-constrained CRFs.}
Constrained decoding answers which legal output should be returned.  Regular
constraints can also be built into a CRF's support; regular-constrained CRFs
are therefore among the closest prior systems \citep{papay2022constraining}.
Our audit keeps the original denominator \(\Z_{\theta}(x)\) and reports
\(\Z_{\theta,\Lang}(x)/\Z_{\theta}(x)\).  A constrained-support CRF normalizes
over legal structures; our diagnostic asks how much legal mass the original
model had before such an intervention.

This is a narrow distinction, not a dominance claim.  If illegal outputs should
receive zero probability by design, a constrained-support CRF is a natural
modeling choice.  Our audit is useful in the complementary situation where the
unconstrained model remains the object under inspection and one wants to know
whether a legal prediction is supported by posterior mass.

\paragraph{Posterior regularization and generalized expectation.}
Posterior regularization can encode posterior constraints and weak structural
supervision \citep{ganchev2010posterior}.  Generalized expectation criteria
also use weak supervision through expectation matching \citep{mann2010ge}.  The
event loss in this paper is a computable use of the audit object, not a new
weak-supervision family.

Posterior regularization is therefore a close conceptual neighbor.  It supplies
an optimization framework for posterior constraints; our main object can be
computed and reported before deciding whether to regularize, project, constrain
decoding, or change model support.  Generalized expectation similarly motivates
care around weak supervision claims: event loss should be described as a use of
posterior event mass, not as a new family of weak-supervision methods.

\paragraph{Semantic loss.}
Semantic loss penalizes low probability assigned to satisfying assignments of a
logical constraint \citep{xu2018semantic}.  The loss
\(-\log \ptheta(\Lang \mid x)\) is semantic-loss-like.  Our emphasis is the CRF
posterior audit object, regular-language diagnostic protocol, and the
distinction from decoded legality and constrained-support normalization.

This concession matters for the paper's scope.  The loss form should not be
presented as unrelated to prior semantic-loss objectives.  The paper instead
asks what this event probability means under a finite CRF posterior and how it
behaves when contrasted with legal decoding, task metrics, and uncertainty
scores.

\paragraph{Constrained optimization and uncertainty.}
Dual decomposition and related constrained optimization methods connect
symbolic constraints with structured prediction objectives
\citep{rush2010dual}.  Our statistic can be evaluated before choosing such an
intervention.  Calibration and uncertainty work ask whether confidence scores
match or rank empirical correctness \citep{platt1999probabilities,guo2017calibration}.
This paper is not a calibration paper: event risk is reported as a
rule-specific posterior-consistency signal, and our diagnostics explicitly show
cases where generic uncertainty baselines are stronger.

\paragraph{Tensor-network motivation.}
Tensor-network language can motivate regular-language events in probabilistic
sequence models, but it is not the main identity of this paper.  The current
results do not establish arbitrary low-rank event transfer, optimized runtime,
or tensor-rank superiority.  Any MPO or rank statement is appendix-only and
conditional on explicit nonnegative factor constructions.

\paragraph{Summary boundary.}
The closest-prior risk is that the paper appears to rename known CRF--automaton
marginal inference and semantic loss.  The manuscript should therefore be read
as a posterior-semantics and auditability paper: the reported object is the
original CRF posterior mass on a regular-language event, and the empirical
claim is about decoded-legality separation, training-signal behavior, and
diagnostic boundaries.

\section{Setup}

Let \(\Y\) be a finite label alphabet and \(T\) a fixed sequence length.  For a
fixed input \(x\), a CRF-style score \(S_{\theta}(x,y)\) assigns a finite real
score to each \(y \in \Y^T\).  The ordinary normalizer is
\[
  \Z_{\theta}(x) = \sum_{y \in \Y^T} \exp S_{\theta}(x,y).
\]
For a regular language \(\Lang \subseteq \Y^*\), define the length-\(T\) event
\(\Lang_T = \Lang \cap \Y^T\) and
\[
  \Z_{\theta,\Lang}(x)
  = \sum_{y \in \Lang_T} \exp S_{\theta}(x,y),
  \qquad
  \ptheta(\Lang \mid x)
  = \frac{\Z_{\theta,\Lang}(x)}{\Z_{\theta}(x)}.
\]

This definition is finite and well-defined for any finite score table.  The
efficient product-transfer computation in the next section additionally assumes
a finite-context local factorization.  An arbitrary finite score table can be
represented by an exponentially large context or explicit table, but that is an
expressivity fallback, not an efficiency claim.

\section{Finite Product Transfer and Event Loss}

The setup in Section 3 defines \(\ptheta(\Lang\mid x)\) for any finite score
table over \(\Y^T\).  The product-transfer computation below requires more
structure.  Assume the score admits a finite-context local factorization.  Let \(C\) be a
finite context state space with initial context \(c_0\), deterministic update
\(U:C \times \Y \to C\), and local potentials
\(\phi_t(x,c_{t-1},y_t)\) such that
\[
  S_{\theta}(x,y) = \sum_{t=1}^T \phi_t(x,c_{t-1},y_t).
\]
Let \(G_t(c,j)=\exp\phi_t(x,c,j)\).  Let a complete DFA for \(\Lang\) be
\((Q,q_{\mathrm{start}},F,\delta)\).  Define the product transfer
\[
M_t^{\Lang}[(c,q),(c',q')]
= \sum_{j \in \Y} G_t(c,j)
  \mathbf{1}[c'=U(c,j)]\mathbf{1}[q'=\delta(q,j)] .
\]
Let \(u_0\) be one at \((c_0,q_{\mathrm{start}})\) and let
\(b_{\Lang}(c,q)=\mathbf{1}[q \in F]\).

The finite-context assumption covers standard first-order linear-chain CRFs and
fixed higher-order Markov contexts.  An arbitrary finite score table can be
represented with a sufficiently large context or table construction, but that
fallback is exponential in general and does not support an efficiency or
scaling claim.  Product-state scaling statements therefore apply only to the
chosen finite-context implementation.

\begin{theorem}
Under the finite setup above, including finite-context local factorization and
a complete deterministic DFA,
\[
  u_0 M_1^{\Lang}\cdots M_T^{\Lang} b_{\Lang}
  = \Z_{\theta,\Lang}(x).
\]
\end{theorem}

\begin{proof}
Fix a label sequence \(y=(y_1,\ldots,y_T)\).  The deterministic context update
and complete DFA transition induce a unique trajectory
\[
  (c_0,q_0),(c_1,q_1),\ldots,(c_T,q_T),
  \qquad q_0=q_{\mathrm{start}},
\]
where \(c_t=U(c_{t-1},y_t)\) and \(q_t=\delta(q_{t-1},y_t)\).  Along this path,
the product of local transfer weights is
\[
  \prod_{t=1}^T G_t(c_{t-1},y_t)
  = \exp\left(\sum_{t=1}^T \phi_t(x,c_{t-1},y_t)\right)
  = \exp S_\theta(x,y).
\]
Conversely, any nonzero term in the expansion of
\(u_0M_1^\Lang\cdots M_T^\Lang b_\Lang\) selects labels that satisfy the same
deterministic context and DFA updates.  If multiple labels induce the same
product transition, the transfer entry sums their weights, and distributivity
of the finite matrix product still assigns each label sequence exactly its own
product of local weights.  The terminal vector \(b_\Lang\) keeps exactly the
trajectories with \(q_T\in F\), which by DFA correctness are exactly
\(\Lang_T\).  Thus the product transfer sums \(\exp S_\theta(x,y)\) over
\(y\in\Lang_T\), which is \(\Z_{\theta,\Lang}(x)\).
\end{proof}

\begin{proposition}
If scores are finite and differentiable and \(\Z_{\theta,\Lang}(x)>0\), then
\[
\nabla_{\theta}[-\log \ptheta(\Lang \mid x)]
= \E_{\ptheta(y\mid x)}[\nabla_{\theta}S_{\theta}(x,y)]
 - \E_{\ptheta(y\mid x,y\in\Lang_T)}[\nabla_{\theta}S_{\theta}(x,y)].
\]
\end{proposition}

\begin{proof}
Because \(\Y^T\) is finite, differentiation commutes with summation.  Therefore
\[
 \nabla_\theta \log \Z_\theta(x)
 = \sum_{y\in\Y^T}\ptheta(y\mid x)\nabla_\theta S_\theta(x,y)
\]
and, when \(\Z_{\theta,\Lang}(x)>0\),
\[
 \nabla_\theta \log \Z_{\theta,\Lang}(x)
 = \sum_{y\in\Lang_T}\ptheta(y\mid x,y\in\Lang_T)
   \nabla_\theta S_\theta(x,y).
\]
Since \(-\log\ptheta(\Lang\mid x)=\log\Z_\theta(x)-\log\Z_{\theta,\Lang}(x)\),
subtracting the two identities gives the result.
\end{proof}

This proposition explains the event loss as a computable training signal.  It
does not imply that task accuracy, F1, calibration, or constrained-decoding
quality improves.

\paragraph{Diagnostic interpretation.}
The scalar \(\ptheta(\Lang\mid x)\) is attached to the original posterior, not
to the selected decoded sequence.  It can therefore be evaluated even when a
constrained decoder returns a legal output.  Low event mass should be
interpreted as posterior inconsistency with a specified rule; whether that
inconsistency predicts task error is an empirical diagnostic question.

\section{Audit and Training Protocol}

We report four quantities separately because they answer different questions.
First, the unconstrained CRF posterior event mass \(\ptheta(\Lang \mid x)\)
measures how much original posterior probability satisfies the rule.  Second,
the legality of the unconstrained Viterbi output records whether the model's
own best sequence is legal.  Third, a constrained-product Viterbi baseline
returns the best legal sequence under the same CRF score and the same rule
language.  Fourth, ordinary task metrics measure agreement with labels.  None
of these quantities subsumes the others.

\subsection{Constrained-Product Decoding Baseline}

The B7 baseline uses the same regular language \(\Lang\) as the audit event and
performs Viterbi decoding over the CRF-context \(\times\) DFA product state.
It returns a legal decoded sequence by construction when a legal path exists.
It does not use \(\ptheta(\Lang \mid x)\) for decoding and should not be read
as a full WFST toolkit or as a replacement for finite-state decoding systems.
Its role is to make decoded legality observable next to original-posterior
event mass.

\subsection{Event Training}

For training experiments, we include a semi-event objective of the form
\[
  \mathcal{L}(\theta)
  = \mathcal{L}_{\mathrm{sup}}(\theta)
    + \lambda[-\log \ptheta(\Lang \mid x)].
\]
Labeled examples use the supervised CRF negative log likelihood.  Rule-labeled
or unlabeled examples can supply event pressure through the second term.  The
gradient proposition in Section 4 explains this as pulling unconstrained
posterior expectations toward event-conditioned posterior expectations.  It is
not a theorem that task accuracy improves.

\subsection{Diagnostic Use}

For diagnostics, we use event risk \(1-\ptheta(\Lang\mid x)\) to rank examples
by posterior inconsistency with the rule.  We report bottom/top quantile gaps,
hidden conflict, and uncertainty baselines when available.  Hidden conflict is
high when legal decoded outputs coexist with low original posterior event mass.
Generic uncertainty baselines are included to prevent over-reading event risk
as a universal uncertainty or calibration score.

\section{Experiments}

The experiments are designed as audit and boundary diagnostics, not as a
benchmark-superiority study.

\paragraph{Hidden posterior conflict.}
WNUT17 diagnostic stress tests show that legal constrained decoding can coexist
with low posterior BIO mass.  These runs are diagnostic only; they do not
support a NER F1 improvement claim.

% Generated by scripts/analysis/generate_latex_tables.py.
% Do not edit numeric values by hand; edit curated CSVs and regenerate.
\begin{table}[t]
\centering
\small
\caption{WNUT17 diagnostic and viability checks.}
\label{tab:table:2:r5:wnut17}
\begin{tabular}{llllllll}
\toprule
block & variant & seeds & P\_event & delta\_P & hidden\_conflict & entity\_F1 & claim\_use \\
\midrule
r5a\_diagnostic\_stress & B0\_unconstrained & 10 & 0.0566 & 0.0000 & 1.0000 & 0.0000 & hidden-conflict diagnostic evidence \\
r5a\_diagnostic\_stress & B4\_semi\_event\_0.1 & 10 & 0.3389 & 0.2822 & 0.9963 & 0.0000 & hidden-conflict diagnostic evidence \\
r5a\_diagnostic\_stress & B5\_rule\_feature\_0.8 & 10 & 0.1608 & 0.1042 & 1.0000 & 0.0000 & hidden-conflict diagnostic evidence \\
r5a\_diagnostic\_stress & B6\_pr\_style\_tau0.7 & 10 & 0.1703 & 0.1137 & 1.0000 & 0.0000 & hidden-conflict diagnostic evidence \\
r5b\_feature\_viability & B0\_unconstrained & 10 & 0.9824 & 0.0000 & 0.0088 & 0.1660 & task viability boundary \\
r5b\_feature\_viability & B4\_semi\_event\_0.1 & 10 & 0.9865 & 0.0041 & 0.0074 & 0.1522 & task viability boundary \\
r5b\_feature\_viability & B5\_rule\_feature\_0.8 & 10 & 0.9864 & 0.0040 & 0.0058 & 0.1645 & task viability boundary \\
r5b\_feature\_viability & B6\_pr\_style\_tau0.7 & 10 & 0.9868 & 0.0044 & 0.0060 & 0.1463 & task viability boundary \\
\bottomrule
\end{tabular}
\end{table}


\paragraph{Event-mass movement.}
Controlled, semi-real, real-source, and public sequence-labeling cases show
that event training can raise posterior event mass.  Task metrics vary, so the
claim is event-mass movement rather than task superiority.

% Generated by scripts/analysis/generate_latex_tables.py.
% Do not edit numeric values by hand; edit curated CSVs and regenerate.
\begin{table}[t]
\centering
\small
\caption{Event-training movement in controlled and field-style probes.}
\label{tab:table:3:event:training:signal}
\begin{tabular}{llllllll}
\toprule
block & family & rows & mean\_delta\_P & positive\_P\_rate & mean\_delta\_legal & mean\_delta\_char & mean\_delta\_exact \\
\midrule
r1\_controlled\_formal & B4 & 8 & 0.1883 & 1 & 0.1639 & 0.0048 & 0.0032 \\
r2\_semi\_real\_formal & B4 & 8 & 0.1584 & 1 & 0.0854 & 0.0188 & 0.0775 \\
r2\_semi\_real\_formal & B5 & 4 & 0.0740 & 1 & 0.0671 & 0.0119 & 0.0449 \\
r2\_semi\_real\_formal & B6 & 4 & 0.0968 & 1 & 0.0772 & 0.0210 & 0.0723 \\
r4\_real\_source\_formal & B4 & 6 & 0.0722 & 1 & 0.0018 & 0.0289 & 0.1127 \\
r4\_real\_source\_formal & B5 & 3 & 0.0314 & 1 & 0.0013 & 0.0187 & 0.0678 \\
r4\_real\_source\_formal & B6 & 3 & 0.0219 & 1 & 0.0009 & 0.0190 & 0.0743 \\
\bottomrule
\end{tabular}
\end{table}


\paragraph{Diagnostic uncertainty boundary.}
In evaluated field-style diagnostics, event risk has positive ranking signal,
but generic uncertainty baselines are stronger in several settings.  We
therefore report event risk as a rule-specific audit statistic, not a universal
uncertainty replacement.

% Generated by scripts/analysis/generate_latex_tables.py.
% Do not edit numeric values by hand; edit curated CSVs and regenerate.
\begin{table}[t]
\centering
\small
\caption{R6a diagnostic ranking and uncertainty boundary.}
\label{tab:table:4:diagnostic:ranking}
\begin{tabular}{lllllllll}
\toprule
source & task & cases & base\_exact\_error & AUROC & AUPRC & Spearman & bottom20\_exact\_error & bottom20\_hidden\_conflict \\
\midrule
all & all & 21000 & 0.7359 & 0.7088 & 0.8470 & 0.2869 & 0.8586 & 0.1033 \\
real\_source & invoice\_6d & 3000 & 0.7203 & 0.7928 & 0.8862 & 0.4949 & 0.9150 & 0 \\
real\_source & invoice\_c6d & 3000 & 0.6547 & 0.7354 & 0.8116 & 0.4347 & 0.8333 & 0.0067 \\
real\_source & stock\_5d & 3000 & 0.7143 & 0.8334 & 0.8979 & 0.5667 & 0.9167 & 0 \\
semi\_real & amount & 3000 & 0.7810 & 0.8046 & 0.9220 & 0.5178 & 0.9600 & 0.0750 \\
semi\_real & date & 3000 & 0.5743 & 0.7220 & 0.7918 & 0.4001 & 0.8817 & 0.4783 \\
semi\_real & dose & 3000 & 0.7867 & 0.8509 & 0.9367 & 0.5296 & 0.9400 & 0.0067 \\
semi\_real & product\_code & 3000 & 0.9200 & 0.8240 & 0.9754 & 0.5165 & 0.9800 & 0.1567 \\
\bottomrule
\end{tabular}
\end{table}

% Generated by scripts/analysis/generate_latex_tables.py.
% Do not edit numeric values by hand; edit curated CSVs and regenerate.
\begin{table}[t]
\centering
\small
\caption{CoNLL2000 public uncertainty boundary.}
\label{tab:table:7:public:uncertainty}
\begin{tabular}{llllll}
\toprule
variant & score & AUROC\_exact & AUPRC\_exact & Spearman\_token\_error & exact\_risk\_gap \\
\midrule
B0\_unconstrained & event\_risk\_1\_minus\_p & 0.7384 & 0.8829 & 0.3117 & 0.4450 \\
B0\_unconstrained & token\_marginal\_entropy & 0.8305 & 0.9388 & 0.4559 & 0.5867 \\
B0\_unconstrained & sequence\_entropy & 0.8339 & 0.9405 & 0.4346 & 0.5817 \\
B0\_unconstrained & viterbi\_margin\_inverse & 0.7860 & 0.9040 & 0.3566 & 0.5167 \\
B0\_unconstrained & max\_sequence\_probability\_inverse & 0.8211 & 0.9334 & 0.4217 & 0.5717 \\
B0\_unconstrained & neg\_log\_viterbi\_probability & 0.8211 & 0.9334 & 0.4217 & 0.5717 \\
B4\_semi\_event\_0.1 & event\_risk\_1\_minus\_p & 0.6895 & 0.8865 & 0.2886 & 0.2650 \\
B4\_semi\_event\_0.1 & token\_marginal\_entropy & 0.8169 & 0.9382 & 0.3972 & 0.5300 \\
B4\_semi\_event\_0.1 & sequence\_entropy & 0.8178 & 0.9394 & 0.3785 & 0.5283 \\
B4\_semi\_event\_0.1 & viterbi\_margin\_inverse & 0.7749 & 0.9063 & 0.3072 & 0.4650 \\
B4\_semi\_event\_0.1 & max\_sequence\_probability\_inverse & 0.8030 & 0.9312 & 0.3590 & 0.5083 \\
B4\_semi\_event\_0.1 & neg\_log\_viterbi\_probability & 0.8030 & 0.9312 & 0.3590 & 0.5083 \\
\bottomrule
\end{tabular}
\end{table}


\paragraph{Public CoNLL2000 case study.}
The public BIO/chunking case provides a multiseed structured-prediction audit.
B4 raises mean posterior BIO event mass and slightly reduces hidden conflict,
while mean token accuracy and span F1 decrease.  This is boundary evidence
against reading event loss as a general accuracy method.

% Generated by scripts/analysis/generate_latex_tables.py.
% Do not edit numeric values by hand; edit curated CSVs and regenerate.
\begin{table}[t]
\centering
\small
\caption{CoNLL2000 public BIO/chunking multiseed case study.}
\label{tab:table:6:public:conll2000}
\begin{tabular}{llllllllll}
\toprule
setting & variant & seeds & P\_BIO\_mean & P\_BIO\_std & hidden\_conflict\_mean & B7\_legal\_mean & token\_acc\_mean & span\_F1\_mean & claim\_use \\
\midrule
formal\_3seed & B0\_unconstrained & 3 & 0.9837 & 0.0011 & 0.0133 & 1.0000 & 0.8731 & 0.7973 & public structured-prediction multiseed audit case; not benchmark superiority \\
formal\_3seed & B4\_semi\_event\_0.1 & 3 & 0.9885 & 0.0063 & 0.0123 & 1.0000 & 0.8697 & 0.7918 & public structured-prediction multiseed audit case; not benchmark superiority \\
\bottomrule
\end{tabular}
\end{table}


\paragraph{Constrained-product baseline and sensitivity.}
B7 reports a faithful constrained-product Viterbi baseline using the same BIO
language.  R7 sensitivity studies show that rule and lambda choice matter,
including a misleading-rule case where event mass rises while task metrics
fall.

% Generated by scripts/analysis/generate_latex_tables.py.
% Do not edit numeric values by hand; edit curated CSVs and regenerate.
\begin{table}[t]
\centering
\small
\caption{B7 constrained-product decoding baseline.}
\label{tab:table:8:b7:constrained:product}
\begin{tabular}{llllll}
\toprule
source\_model & decoded\_legal\_rate & token\_accuracy & entity\_F1 & uses\_event\_training & uses\_event\_mass\_for\_decoding \\
\midrule
B0\_unconstrained & 1.0000 & 0.8841 & 0.1486 & False & False \\
B4\_semi\_event\_0.1 & 1.0000 & 0.8838 & 0.1688 & True & False \\
\bottomrule
\end{tabular}
\end{table}

% Generated by scripts/analysis/generate_latex_tables.py.
% Do not edit numeric values by hand; edit curated CSVs and regenerate.
\begin{table}[t]
\centering
\small
\caption{R7 lambda/rule sensitivity and boundary cases.}
\label{tab:table:9:r7:sensitivity}
\begin{tabular}{lllllllll}
\toprule
task & variant & runs & P\_event & delta\_P\_event & delta\_legal\_rate & delta\_char\_acc & delta\_exact\_acc & boundary \\
\midrule
date & B0\_unconstrained & 10 & 0.7136 & 0.0000 & 0.0000 & 0.0000 & 0.0000 & ordinary \\
date & B4\_semi\_event\_1.0 & 10 & 0.9738 & 0.2603 & 0.1948 & 0.0306 & 0.1996 & ordinary \\
product\_code & B0\_unconstrained & 10 & 0.7765 & 0.0000 & 0.0000 & 0.0000 & 0.0000 & ordinary \\
product\_code & B4\_semi\_event\_1.0 & 10 & 0.9844 & 0.2079 & 0.0328 & 0.0097 & 0.0248 & ordinary \\
product\_code\_swapped\_rule & B0\_unconstrained & 10 & 0.0000 & 0.0000 & 0.0000 & 0.0000 & 0.0000 & event/task tradeoff \\
product\_code\_swapped\_rule & B4\_semi\_event\_0.1 & 10 & 0.6083 & 0.6083 & 0.6752 & -0.4352 & -0.0668 & event/task tradeoff \\
product\_code\_swapped\_rule & B4\_semi\_event\_1.0 & 10 & 0.9486 & 0.9485 & 0.9716 & -0.5524 & -0.0816 & event/task tradeoff \\
stock\_like\_digits & B0\_unconstrained & 10 & 0.9192 & 0.0000 & 0.0000 & 0.0000 & 0.0000 & legal-rate-not-useful \\
stock\_like\_digits & B4\_semi\_event\_1.0 & 10 & 0.9961 & 0.0769 & 0.0000 & 0.0331 & 0.0868 & legal-rate-not-useful \\
\bottomrule
\end{tabular}
\end{table}


\section{Limitations}

This paper does not introduce a new product-automaton inference algorithm.
Product automaton marginal inference is the known computational neighborhood.
The method is not a replacement for WFST methods, constrained decoding,
regular-constrained CRFs, posterior regularization, or semantic loss.

The empirical evidence is not benchmark-superiority evidence.  WNUT17
diagnostic stress has zero entity F1 in the hidden-conflict setting, the public
CoNLL2000 multiseed case shows event-mass movement with lower mean task metrics
for B4, and uncertainty baselines can be stronger than event risk.  The tensor
or MPO material is appendix-only motivation and does not establish arbitrary
low-rank event transfers or optimized runtime.

Final submission requires a fresh proof audit, citation-context audit, and
paper-claim audit on the completed LaTeX and BibTeX package.

\section{Conclusion}

Posterior regular-language event mass provides a reportable audit statistic for
CRF posteriors.  It makes explicit a distinction that constrained decoding can
hide: a legal decoded output need not mean the original posterior assigns high
mass to legal structures.  Event loss offers a computable training signal, but
the experiments support a conservative reading: event mass is useful as
posterior semantics and rule-specific audit evidence, with clear empirical
boundaries.


\bibliographystyle{plainnat}
\bibliography{references}

\appendix
\section{Reproducibility Notes}

The paper tables are generated from curated CSV files under
\texttt{experiments/results/paper\_tables/}.  The LaTeX fragments under
\texttt{paper/tables/} are generated by
\texttt{scripts/analysis/generate\_latex\_tables.py}.  Raw run bundles are not
tracked in Git; curated audit files record provenance, commit hashes, configs,
return codes, durations, and claim boundaries.

Before submission, regenerate the table fragments and rerun:
\[
\texttt{python scripts/analysis/verify\_jmlr\_claim\_evidence.py}.
\]


\end{document}
